\section{Objectives}
\label{sec:objectives}

The aim is to design and develop a gamified web-security training platform that lets
students practise penetration-testing techniques on realistic vulnerable web
applications, while providing interactive learning, performance evaluation, and
instructor analytics. This breaks into six objectives. Each closes a gap
(Table~\ref{tab:gaps}) and carries a measurable success criterion. The ``adaptive''
character of the platform is delivered by O3--O5 acting together: scoring produces a
performance record (O3), a simple model turns it into a prediction and personalised
recommendation (O4), and the instructor sees the result (O5).

\begin{enumerate}[label=\textbf{O\arabic*.},leftmargin=2.6em,itemsep=0.45em]

\item \textbf{Vulnerable web application (closes G1).}
Develop an intentionally vulnerable web application incorporating common OWASP
Top~10 vulnerabilities, delivered as challenges inside a single classroom platform.
\emph{Success:} OWASP Top~10 vulnerability families (for example SQL injection,
broken access control, cross-site scripting), each with challenges that are
exploitable end to end in a contained environment.

\item \textbf{Gamification module (closes G2).}
Design a gamification module that includes points, badges, levels, achievements, and
leaderboards, deliberately designed against the documented negative effects of
gamification (for example, a peer-banded leaderboard rather than one absolute
ranking, and optional rather than forced competition).
\emph{Success:} all five game elements built, with at least one design choice
per element justified against the negative-effects literature.

\item \textbf{Automated scoring (closes G3).}
Develop an automated scoring mechanism that evaluates learner performance from
challenge completion, accuracy, and time taken, recording each attempt as the
per-learner performance record the adaptive layer will use.
\emph{Success:} every attempt is scored automatically with no manual grading, and
the score is a function of completion, accuracy, and time.

\item \textbf{Adaptive hints, resources, and recommendation (closes G4).}
Provide hints and learning resources to assist students during exercises, and use a
simple, interpretable machine-learning model over the performance record (O3) to
predict a learner's likely difficulty and recommend a suitably difficult next
challenge. Assignments are customised because every learner is different.
\emph{Success:} a working hint/resource layer that unlocks on struggle, plus a
model whose next-challenge recommendation is validated against held-out learner
data.

\item \textbf{Instructor dashboard (closes G5).}
Develop an instructor dashboard to monitor learner progress, challenge completion,
and overall performance, surfacing the model's prediction of which learners are at
risk so that support arrives early.
\emph{Success:} a dashboard showing per-learner and cohort progress, completion, and
scores, with at-risk learners flagged.

\item \textbf{Evaluation of effectiveness (closes G6).}
Evaluate the effectiveness of the gamified platform in improving learner engagement
and practical cybersecurity skills.
\emph{Success:} a pilot with a student cohort reporting a pre/post change in
practical skill and a measured engagement outcome, as set out in
Section~\ref{sec:methodology}.

\end{enumerate}
